%7. Willingness to contribute: majority willing to contribute a few percent

%===============================================================================================
% - we are not interested in WTP because what matters is willingness to support policy (which makes the richest pay, not everyone for fairness reasons) and willingness to contribute in case it's a policy and the others are also put to contribution
%===============================================================================================

This section examines which international policies citizens are willing to support and how their support varies with policy design, including contribution rules, country participation, effectiveness, fairness, and redistribution. Accordingly, the relevant outcome is policy support rather than the amount an individual is willing to pay. \textit{Willingness to pay} captures an individual's trade-off between a monetary sacrifice and the achievement of a given objective, whereas \textit{willingness to support} refers more broadly to the acceptance of a collective policy. Individuals may support a policy without being net contributors. For many of the international policies examined thus far, contribution rules protect lower-income households and assign a larger share of the costs to higher-income individuals, wealthier taxpayers, and countries with greater financial capacity. Public support therefore depends partly on considerations of fairness, especially the progressivity of contributions and the distribution of the fiscal burden within and across countries. \textit{Willingness to support} is thus the more appropriate concept for international policies, since contributions vary with ability to pay and cannot be represented by a uniform \textit{willingness to pay}. The remainder of this section therefore examines how far citizens are prepared to support such collective arrangements, moving from a uniform contribution of 1\% of income to the use of domestic taxation for global problems and, finally, to a self-designed global redistribution that may reduce respondents' own income.

%===============================================================================================
% - 1% for climate action (Andre et al)
%===============================================================================================

%LB: J'ai aussi pris en compte les commentaires que tu m'avais laissé dans la V1 du draft.
Drawing on the 2021--2022 Gallup World Poll, \citet{andre_globally_2024} analyze responses from representative samples of people aged 15 and older across 125 countries. The survey asks respondents whether they would contribute 1\% of their household income each month to fight global warming. In the population-weighted sample, 69\% accept this contribution. An additional 6\% would contribute a smaller amount, while 26\% would contribute nothing. Willingness to contribute the full 1\% exceeds 50\% in 114 of the 125 countries and ranges from 30\% to 93\%. The eleven countries below the majority threshold are Japan (49\%), Canada (49\%), the United States (48\%), the United Kingdom (48\%), Pakistan (47\%), New Zealand (46\%), Kazakhstan (45\%), Russia (41\%), Lithuania (41\%), Israel (37\%), and Egypt (30\%) \citep[Figure 1 \& Supplementary Table S4]{andre_globally_2024}. 

Misperceptions about collective participation may help explain some of the reluctance to contribute. Although 69\% of respondents are willing to contribute, they estimate that only 43\% of their fellow citizens would do so, a gap of 26 percentage points. Actual willingness exceeds perceived willingness in every surveyed country. Moreover, respondents' beliefs about others' willingness correlate strongly with their own willingness to contribute \citep[Figure 3 \& Supplementary Table S11]{andre_globally_2024}.

%/LB: p.10-11 du rapport de la WB (2009) => argument global governance (section 5, Ghassim) via 'Evaluation Government Action'.
An earlier World Bank survey covering 15 countries yields similar results. Respondents first indicated whether they would accept higher prices for energy and other products to fight climate change. Within each country, all respondents were presented with the same nominal monthly amount, calculated as 1\% of that country's annual GDP per capita and prorated monthly; those who rejected it then considered a second fixed amount equivalent to 0.5\% of GDP per capita. Because the amount did not vary with individual income, this design is regressive within countries and may therefore yield conservative estimates of support relative to a proportional contribution. Willingness to accept the initial amount ranged from 11\% in Russia to 68\% in China. Acceptance of the full amount was highest in China (68\%), Vietnam (59\%), Japan (53\%), and Iran and Mexico (51\% each), whereas it was lowest in Russia (11\%) and Bangladesh (32\%). The pattern thus does not follow a simple divide between richer and poorer countries. When the lower amount is included, support reaches a majority in 14 of the 15 countries and is highest in Vietnam (85\%) and China (82\%), while Russia is the only country without majority support (25\%). Overall, 46\% accepted the full amount, 63\% accepted either amount, and 33\% rejected both \citep{world_bank_public_2009}. Because all respondents within a country were asked to pay the same absolute amount regardless of their income, the proposal is regressive, imposing a proportionally greater burden on poorer individuals and thus probably yielding a conservative estimate of support.
 % AF: so they propose the same absolute amount to everyone instead of a relative (%) amount? If so, we should flag this as probably yielding conservative results given the regressivity of the proposal. Also, give more info on the country specific attitudes (which types of country have low vs. high support)
%LB: respondents face the same relative amount (1\% of their country GDP per capita), so there's indeed an issue of regressivity.

%===============================================================================================
% - Most agree own taxes should pay for global issues (Fabre 26)
%===============================================================================================

%LB: j'ai ajouté un contraste avec ce qu'on avait dit avant dans la section HICs financing LICs. Ça fait un rappel avec le premier paragraphe de cette section.
\citet{fabre_public_2026} asked respondents whether they agreed that their ``taxes should go towards solving global problems''. Excluding neutral responses, 59\% agree and 41\% disagree \citep[Figures S59--S60]{fabre_public_2026}.
Earlier results show substantially higher acceptance when the contribution falls on the richest: policies requiring contributions from the global top 1\% and top 3\% receive acceptance rates of 69\% and 64\%, respectively. This contrast suggests that support depends on how the policy distributes the contribution burden.

%===============================================================================================
% - Almost half choose a custom redistr that makes them lose (Fabre 26)
%===============================================================================================

%LB: ATTENTION c'est quelque chose que j'ai déjà cité dans la section 4, donc il faudra l'enlever et la remettre ici (j'imagine). J'ai fait un copier-coller de la section 4.
\citet{fabre_public_2026} then asked respondents to construct their preferred global distribution of after-tax income by choosing the shares of winners and losers and the redistribution from richer to poorer individuals. Respondents see both the current and proposed distributions: ``The current distribution is in red, and your custom one is in green'' and could preserve the status quo. The interactive exercise is available at \url{http://preferences-pol.fr/custom_global_redistr.html}, with an explanatory video at \url{https://youtu.be/gSfsQwczT7w}. The median preferred parameters are to make 49\% of the world population better off and 18\% worse off, transferring 5\% of world income from richer to poorer individuals and raising the minimum monthly income to \$287 at purchasing power parity \citep[Figure S49]{fabre_public_2026}. Almost half of respondents (48\%) select a redistribution that would reduce their own income, compared with 9\% who would personally gain. Only 10\% choose to retain the current distribution and not reduce inequality. The average preferred redistribution would likewise transfer 5.4\% of world income, in this case from the top 27\% to the bottom 73\%, and establish a minimum monthly income of \$247 \citep[Figure S50]{fabre_public_2026}. The average and median proposals would reduce the global Gini coefficient from 0.67 to 0.59 and 0.58, respectively. 







